\documentclass[../main.tex]{subfiles}
\begin{document}

\chapter{フォールトトレランス}
\lessoninfo{
  video={Lesson 17，動画 1--32},
  textbook={H\&P \cite{hennessy2017} §1.7，§D.2--D.3},
  keywords={ディペンダビリティ，フォールト，エラー，故障，信頼性，可用性，MTTF，MTTR，フォールトトレランス，チェックポインティング，N モジュール冗長，パリティ，ECC，RAID},
}

第16章では，データを永続的に保存する装置を扱った．ストレージ装置を含め，
計算機のあらゆる構成要素は，いずれ意図された振る舞いから逸脱する．構成
要素が1つでもそうなるたびに誤った結果を返したりデータを失ったりする
システムは，信頼して使うことができない．この章では，そうした問題を論じる
ための用語---フォールト，エラー，故障，および信頼性と可用性---を導入し，
実際に起こるフォールトを分類し，それらを許容する一般的な技法を示す．後半
ではこれらの考え方をディスクに適用する．独立ディスク冗長アレイ（RAID）は
複数のディスクを1台にまとめたものであり，その冗長な構成により，ディスクが
丸ごと故障してもデータは失われず，アクセスも可能なままに保たれる．

\section{ディペンダビリティ}
\label{sec:l17-dependability}

計算機が有用であるのは，その利用者が，計算機の提供するサービスを信頼
できる場合に限られる．\source{Lesson 17，動画 1--2；H\&P §1.7；
\cite{avizienis2004}．}これを厳密にするために，システムの振る舞いに関する
2つの記述を区別する．\emph{仕様サービス}とは，仕様に従えばシステムが持つ
べき振る舞いであり，\emph{提供サービス}とは，システムが実際に示す振る舞い
である．

\begin{definition}[ディペンダビリティ]
  システムの\term[でぃぺんだびりてぃ]{ディペンダビリティ}[dependability]と
  は，そのサービスの提供についてシステムに依拠することを正当化する，提供
  サービスの品質である．
\end{definition}

システムは構成要素，すなわち\emph{モジュール}から作られる．プロセッサ，
メモリモジュール，ディスク，さらに細かく見ればそれらの内部の回路や
プログラムの関数である．各モジュールはそれ自身の仕様上の振る舞いを持つ．
ディペンダビリティの問題は，あるモジュールがその仕様から逸脱したときに
始まる．そのような逸脱がシステム全体の逸脱に変わりうること，しかし必ず
そうなるわけではないことを，次節で述べる．

\section{フォールト・エラー・故障}
\label{sec:l17-fault-error-failure}

1つのモジュールの内部で始まった問題は，システムの利用者が観測できるように
なるまでに最大3つの段階を経る．\source{Lesson 17，動画 3--5；H\&P §1.7，
§D.3．}段階に別々の名前が与えられているのは，ディペンダブルなシステムが，
段階の間の遷移のいずれかで問題を食い止めようとするからである．

\begin{definition}[フォールト・エラー・故障]
  \label{def:l17-fault-error-failure}
  \term[ふぉーると]{フォールト}[fault]とは，モジュールが仕様上の振る舞い
  から逸脱することである．欠陥のあるトランジスタや，プログラムの誤りが
  その例である．

  \term[えらー]{エラー}[error]とは，システム内部の実際の振る舞い---その
  内部状態---が仕様上の振る舞いと異なっている状態である．変数に誤った値が
  格納されている場合がその例である．

  \term[こしょう]{故障}[failure]とは，システムの提供サービスが仕様サービス
  から逸脱することである．
\end{definition}

フォールトは，何の影響も及ぼさないまま長期間存在しうる．フォールトのある
モジュールが，その逸脱を露わにする形で使われるまでは，フォールトは
\emph{潜在}エラーである．そのような使われ方をすると，フォールトは
\emph{活性化}され，\emph{顕在}エラーとなる．内部状態がこの時点で誤った
ものになる．顕在エラーが故障を引き起こすのは，それがシステムの提供する
サービスまで伝播した場合だけである．すべてのフォールトが活性化されるわけ
ではなく，すべてのエラーが故障になるわけでもない
（\cref{fig:l17-fault-chain}）．

\begin{example}
  \label{ex:l17-fault-chain}
  2つの整数を掛ける関数にプログラムの誤りがある．引数が 6 と 9 の場合を
  除けば正しい積を返すが，この場合だけ 54 ではなく 45 を返す．この誤りは
  関数が書かれた時点からフォールトであり，どの呼び出し元も 6 と 9 を掛け
  ない限り潜在したままである．プログラムがこの引数で関数を呼び出し，結果を
  格納すると，フォールトは活性化され，プログラムの状態はエラーを含む．
  その結果が予約システムの確保する座席数であれば，54 席ではなく 45 席が
  確保される．エラーは提供サービスに到達し，故障が起きたことになる．一方，
  プログラムが積の正負だけを判定するのであれば，誤った値でも正しい値と
  同じ判断に至り，故障は起きない．
\end{example}

\begin{figure}[htbp]
  \centering
  % Fault -> error -> failure.  An error that leaves a module is a failure of
% that module and, at the same time, a fault of the enclosing system; an error
% that reaches the system's delivered service is a failure of the system.
\begin{tikzpicture}[x=1cm,y=1cm, font=\footnotesize\sffamily, >={Stealth[length=1.8mm]},
  st/.style={draw, rounded corners=2pt, fill=white, minimum height=0.65cm,
    minimum width=1.2cm, inner sep=2pt},
  lab/.style={font=\scriptsize, text=ln-muted, align=center}]
  % boundaries
  \draw[ln-rule, fill=ln-box] (0,0) rectangle (9.0,3.2);
  \node[anchor=north west, font=\scriptsize\sffamily\bfseries, text=ln-accent]
    at (0.05,3.15) {\iflangja{システム}{System}};
  \draw[ln-rule, fill=ln-accent!10] (0.25,1.05) rectangle (4.8,2.75);
  \node[anchor=north west, font=\scriptsize\sffamily\bfseries, text=ln-accent]
    at (0.3,2.7) {\iflangja{モジュール}{Module}};
  % chain
  \node[st] (f)  at (1.1,1.75) {\iflangja{フォールト}{fault}};
  \node[st] (e1) at (3.85,1.75) {\iflangja{エラー}{error}};
  \node[st] (e2) at (7.0,1.75) {\iflangja{エラー}{error}};
  \node[st, draw=ln-caution, text=ln-caution] (fl) at (10.6,1.75) {\iflangja{故障}{failure}};
  \draw[->] (f) -- node[lab, above=1pt] {\iflangja{活性化}{activated}} (e1);
  \draw[->] (e1) -- node[lab, above=1pt, pos=0.55] {\iflangja{伝播}{propagates}} (e2);
  \draw[->] (e2) -- node[lab, above=1pt, pos=0.28] {\iflangja{伝播}{propagates}} (fl);
  % boundary crossings
  \fill[ln-accent] (4.8,1.75) circle (1.6pt);
  \fill[ln-caution] (9.0,1.75) circle (1.6pt);
  \node[lab, anchor=north] at (4.8,0.95)
    {\iflangja{モジュールの故障＝\\システムのフォールト}{module failure =\\fault of the system}};
  \node[lab, anchor=north] at (10.6,1.3)
    {\iflangja{提供サービスが\\仕様から逸脱}{delivered service\\deviates from\\specified service}};
  \node[lab, anchor=north] at (1.1,1.38) {\iflangja{（潜在）}{(latent)}};
\end{tikzpicture}

  \caption{フォールトは活性化されるとエラーになる．モジュールの外に出た
    エラーは，そのモジュールの故障であると同時に，そのモジュールを含む
    システムのフォールトである．提供サービスに到達したエラーはシステムの
    故障である．}
  \label{fig:l17-fault-chain}
\end{figure}

この3つの用語は，どの境界を考えるかに相対的である．モジュールの故障は，
そのモジュールを含むシステムから見れば，そのシステムのフォールトである．
ヘッドがクラッシュしたディスクは，ディスクとしては故障している．しかし，
データの冗長な複製を保持するディスクアレイ（\cref{sec:l17-raid}）にとって，
同じ事象はフォールトであり，アレイの故障につながるとは限らない．

\section{信頼性と可用性}
\label{sec:l17-reliability}

システムは，その仕様サービスに関して，任意の時点で2つの状態のいずれかに
ある．\source{Lesson 17，動画 6--7；H\&P §1.7．}\emph{サービス遂行}の
状態では提供サービスが仕様サービスと一致しており，\emph{サービス中断}の
状態では一致していない．故障はシステムを第1の状態から第2の状態へ移し，
復旧（修復）はシステムを元の状態に戻す
（\cref{fig:l17-service-timeline}）．それぞれの状態で過ごす時間から，
ディペンダビリティの異なる2つの尺度が得られる．\margin{このモデルに第3の
状態はない．遅いが正しいサービスや，一部の利用者にだけ中断しているサービス
も，どちらかに数えるほかない．}

\begin{figure}[htbp]
  \centering
  % Two service states and a timeline alternating between them.  Up-times are
% times to failure, down-times are times to repair.
\begin{tikzpicture}[x=1cm,y=1cm, font=\footnotesize\sffamily, >={Stealth[length=1.8mm]},
  st/.style={draw, rounded corners=3pt, minimum height=0.65cm, inner sep=3pt},
  lab/.style={font=\scriptsize, text=ln-muted, align=center}]
  % states
  \node[st, fill=ln-accent!15] (up) at (2.6,2.3) {\iflangja{サービス遂行}{service accomplishment}};
  \node[st, fill=ln-caution!15] (down) at (8.4,2.3) {\iflangja{サービス中断}{service interruption}};
  \draw[->] (up) to[bend left=12] node[lab, above] {\iflangja{故障}{failure}} (down);
  \draw[->] (down) to[bend left=12] node[lab, below] {\iflangja{修復}{restoration}} (up);
  % timeline
  \foreach \a/\b/\n in {0/3.1/1, 3.6/6.4/2, 7.2/10.6/3} {
    \fill[ln-accent!15] (\a,0) rectangle (\b,0.45);
    \draw[|<->|, ln-muted] (\a,0.7) -- node[above, font=\scriptsize] {$U_\n$} (\b,0.7);
  }
  \foreach \a/\b/\n in {3.1/3.6/1, 6.4/7.2/2, 10.6/11.1/3} {
    \fill[ln-caution!40] (\a,0) rectangle (\b,0.45);
    \node[font=\scriptsize, below] at ({(\a+\b)/2},-0.05) {$D_\n$};
  }
  \draw (0,0) rectangle (11.1,0.45);
  \draw[->] (0,-0.55) -- (11.4,-0.55) node[pos=1, below left, lab] {\iflangja{時間}{time}};
  \node[lab, anchor=north] at (5.55,-0.75)
    {MTTF $=$ \iflangja{$U_i$ の平均}{mean of $U_i$}\qquad
     MTTR $=$ \iflangja{$D_i$ の平均}{mean of $D_i$}};
\end{tikzpicture}

  \caption{2つのサービス状態と，それらを交互に繰り返す履歴．$U_i$ が
    サービス遂行の期間，$D_i$ がサービス中断の期間である．}
  \label{fig:l17-service-timeline}
\end{figure}

\begin{definition}[信頼性]
  システムの\term[しんらいせい]{信頼性}[reliability]とは，サービスが
  \emph{連続して}遂行されることの尺度である．通常は
  \term[へいきんこしょうじかん]{平均故障時間}[mean time to failure]%
  \index{MTTF|see{平均故障時間}}（MTTF）で定量化される．MTTF とは，サービス
  遂行の期間の平均の長さ，すなわちサービス開始から次の故障までの平均時間で
  ある．
\end{definition}

\begin{definition}[可用性]
  システムの\term[かようせい]{可用性}[availability]とは，サービス遂行に
  費やす時間の割合である．サービス中断の期間の平均の長さを
  \term[へいきんしゅうふくじかん]{平均修復時間}[mean time to repair]%
  \index{MTTR|see{平均修復時間}}（MTTR）とすると，
  \begin{equation}
    \text{可用性} \;=\; \frac{\mathrm{MTTF}}{\mathrm{MTTF}+\mathrm{MTTR}} .
    \label{eq:l17-availability}
  \end{equation}
\end{definition}

2つの平均の和は，故障と修復の1周期の平均の長さであり，
\term[へいきんこしょうかんかく]{平均故障間隔}[mean time between failures]%
\index{MTBF|see{平均故障間隔}}と呼ばれる．
\begin{equation}
  \mathrm{MTBF} \;=\; \mathrm{MTTF} + \mathrm{MTTR} .
  \label{eq:l17-mtbf}
\end{equation}
修復がサービスの期間に比べて短ければ，MTBF と MTTF はほぼ等しい．
\caution{ディスクやサーバのデータシートが「MTBF」として示す値は，ここで
MTTF と定義した量である．修復時間は含まれていない．}

\begin{example}
  \label{ex:l17-availability}
  あるサーバが \SI{1200}{\hour} 稼働し，\SI{6}{\hour} 停止し，
  \SI{900}{\hour} 稼働し，\SI{10}{\hour} 停止し，\SI{1500}{\hour} 稼働し，
  \SI{8}{\hour} 停止した．MTTF は $(1200+900+1500)/3 = \SI{1200}{\hour}$，
  MTTR は $(6+10+8)/3 = \SI{8}{\hour}$，MTBF は \SI{1208}{\hour} であり，
  可用性は $1200/1208 \approx \SI{99.34}{\percent}$ である．\tip{可用性は
  「ナイン」の数で語られる．1 年の \SI{8766}{\hour} のうち，
  \SI{99.9}{\percent} は \SI{8.8}{\hour}，\SI{99.99}{\percent} は
  \SI{53}{\minute}，\SI{99.999}{\percent} は \SI{5}{\minute} の中断を許す．}%
  故障頻度が
  10 倍（MTTF $=\SI{120}{\hour}$）でありながら平均 \SI{0.8}{\hour} で復旧
  する別のサーバは，可用性は同じだが MTTF は 10 分の 1 であり，信頼性は
  はるかに低い．
\end{example}

\section{フォールトの種類}
\label{sec:l17-kinds}

フォールトは2通りに分類される．原因による分類と，どれだけ長く続くかに
よる分類である．\source{Lesson 17，動画 8--9；H\&P §D.3．}

\subsection{原因による分類}

\begin{description}
  \item[ハードウェア] \term[はーどうぇあふぉーると]{ハードウェアフォールト}[hardware fault]%
    は，ハードウェアが設計どおりに動作しなくなったときに生じる．スイッチ
    しなくなったトランジスタや，プラッタに接触した磁気ヘッドがその例で
    ある．
  \item[設計] \term[せっけいふぉーると]{設計フォールト}[design fault]とは，
    設計そのものの誤りである．ソフトウェアのバグや，初期の Intel Pentium
    の除算バグのようなハードウェア設計の誤りがこれにあたる．後者は，欠陥の
    ある設計から作られたすべてのチップに存在した．
  \item[運用] \term[うんようふぉーると]{運用フォールト}[operation fault]と
    は，システムの運用者や利用者が犯す誤りである．誤ったファイルを削除
    する，誤ったケーブルを抜く，などがその例である．
  \item[環境] \term[かんきょうふぉーると]{環境フォールト}[environmental fault]%
    は，システムの外部に由来する．火災，浸水，停電，極端な温度，意図的な
    破壊行為などである．
\end{description}

\subsection{継続時間による分類}

\begin{description}
  \item[永久] \term[えいきゅうふぉーると]{永久フォールト}[permanent fault]%
    は，いったん起きると，フォールトのある部品が修理または交換されるまで
    続く．過熱で損傷したチップがその例である．
  \item[間欠] \term[かんけつふぉーると]{間欠フォールト}[intermittent fault]%
    は，しばらく続いた後に消え，また再発する．緩んだコネクタや，定格を
    超えるクロック周波数で動作し，高温になったときだけ誤った計算を行う
    プロセッサがその例である．
  \item[過渡] \term[かとふぉーると]{過渡フォールト}[transient fault]は
    一度だけ起き，ひとりでに消える．高エネルギー粒子によって反転した
    メモリのビットがその例であり，そのビットを書き直せば正しい動作に戻る．
    \caution{メモリの誤りの多くはこの種のものではない．Google のサーバの
    実地調査では，年間 \SI{8}{\percent} 超の DIMM が訂正可能な誤りに見舞われ，
    その多くは粒子の衝突ではなく素子側の恒久的な欠陥であった
    \cite{schroeder2009}．}
\end{description}

2つの分類は，フォールトの異なる性質を記述するものである．例えばハード
ウェアフォールトは，永久（焼き切れたチップ），間欠（温まったときだけ接触が
失われるはんだ割れ），過渡（反転したメモリのビット）のいずれにもなりうる．
一方で設計フォールトは，設計が修正されるまで設計のすべての複製に存在する
ので，たまにしか活性化されないとしても永久フォールトである．\margin{Pentium
の除算器では，\num{2048} エントリの参照表のうち 5 エントリが誤った値を
保持していた．チップ交換に要した費用は 4 億 7500 万ドルである．}

\section{信頼性と可用性の向上}
\label{sec:l17-improving}

システムをよりディペンダブルにする方法は3つある．\source{Lesson 17，
動画 10；H\&P §1.7．}

\begin{description}
  \item[フォールトを避ける] そもそもフォールトが起きないようにすることを
    \term[ふぉーるとかいひ]{フォールト回避}[fault avoidance]という．故障率
    の低い部品を使う，部品を低温に保つ，機械に液体を近づけない，などで
    ある．
  \item[フォールトを許容する] フォールトが故障になることを防ぐことを
    \term[ふぉーるととれらんす]{フォールトトレランス}[fault tolerance]と
    いう．一部のモジュールにフォールトがあっても，システムは仕様サービスを
    提供し続ける．フォールトトレランスは冗長性の上に成り立つ．メモリの
    誤り訂正符号（\cref{sec:l17-codes}）や冗長なディスク
    （\cref{sec:l17-raid}）がその例である．
  \item[修復を速くする] 故障からサービス復旧までの時間を短くする．予備
    部品を現地に置いておく，などである．
\end{description}

\begin{keypoint}
  フォールト回避とフォールトトレランスは故障の回数を減らすので，MTTF を
  高め，それによって信頼性と可用性の両方を高める．修復の高速化は MTTR を
  下げるだけである．モジュールが故障するたびに停止するシステムでは，
  可用性は改善するが信頼性は改善しない．
\end{keypoint}

\section{フォールトトレランスの技法}
\label{sec:l17-techniques}

汎用的なフォールトトレランスは2つの要素を組み合わせる．エラーを
\emph{検出}する手段と，エラーから\emph{回復}する手段である．
\source{Lesson 17，動画 11--13．}

\subsection{チェックポインティング}

システムの完全な状態を一定間隔で保存し，後でその保存状態から実行を再開
できるようにすることを\term[ちぇっくぽいんてぃんぐ]{チェックポインティング}[checkpointing]%
という．エラーが検出されると，システムは直近のチェックポイントを復元し，
そこから実行をやり直す．\margin{プロセッサは数命令ごとにこれを行っている．
第8章の分岐予測ミス回復は，レジスタマップのチェックポイントを復元して実行
し直す．違うのは粒度だけである．}フォールトが過渡的であれば，やり直しは
成功する見込みが高く，間欠的であっても多くの場合は成功する．永久フォールト
は再び活性化されるので，この方法では許容できない．\caution{「過渡」が表す
のはフォールトであって，あとに残るエラーではない．誤った値は何かが上書き
するまで残る．やり直しが効くのは，チェックポイントがエラーより前にあるから
にすぎない．}回復にも時間がかかる．
チェックポイントの復元と失われた作業のやり直しが速ければ，利用者は遅延を
感じるだけで済むが，時間がかかりすぎれば，回復そのものがサービス中断と
みなされる．

\subsection{冗長化}

\begin{definition}[$N$ モジュール冗長]
  \term[N もじゅーるじょうちょう]{$N$ モジュール冗長}[N-module redundancy]%
  では，同一の $N$ 個のモジュールが同じ入力に対して同じ計算を行い，その
  出力を比較するか多数決にかける（\cref{fig:l17-nmr}）．
\end{definition}

\begin{figure}[htbp]
  \centering
  % N-module redundancy: (a) two modules and a comparator detect an error;
% (b) three modules and a majority voter mask the error of one module.
\begin{tikzpicture}[x=1cm,y=1cm, font=\footnotesize\sffamily, >={Stealth[length=1.8mm]},
  mod/.style={draw, fill=ln-box, minimum width=1.1cm, minimum height=0.55cm, inner sep=2pt},
  lab/.style={font=\scriptsize, text=ln-muted, align=center}]
  % ---------------- (a) dual
  \node[anchor=west, font=\footnotesize\bfseries] at (-0.5,1.55) {(a) $N=2$};
  \draw (-0.5,0) node[left, lab] {\iflangja{入力}{in}} -- (0.3,0);
  \fill (0.3,0) circle (1.2pt);
  \node[mod] (a1) at (1.3,0.55) {M1};
  \node[mod] (a2) at (1.3,-0.55) {M2};
  \draw[->] (0.3,0) |- (a1.west);
  \draw[->] (0.3,0) |- (a2.west);
  \node[draw, circle, fill=ln-accent!15, minimum size=0.7cm, inner sep=0pt] (cmp) at (3.1,0) {$=$};
  \draw[->] (a1.east) -| (cmp.north);
  \draw[->] (a2.east) -| (cmp.south);
  \draw[->] (cmp.east) -- (4.0,0) node[right, lab] {\iflangja{出力}{out}};
  \node[lab, anchor=north] at (3.1,-0.9)
    {\iflangja{一致: 結果は有効\\不一致: エラー検出}{equal: result valid\\differ: error detected}};
  % ---------------- (b) triple
  \begin{scope}[xshift=6.0cm]
  \node[anchor=west, font=\footnotesize\bfseries] at (-0.5,1.55) {(b) $N=3$};
  \draw (-0.5,0) node[left, lab] {\iflangja{入力}{in}} -- (0.3,0);
  \fill (0.3,0) circle (1.2pt);
  \node[mod] (b1) at (1.3,0.8) {M1};
  \node[mod] (b2) at (1.3,0) {M2};
  \node[mod] (b3) at (1.3,-0.8) {M3};
  \draw[->] (0.3,0) |- (b1.west);
  \draw[->] (0.3,0) -- (b2.west);
  \draw[->] (0.3,0) |- (b3.west);
  \node[draw, fill=ln-accent!15, minimum width=0.9cm, minimum height=2.1cm, align=center, inner sep=1pt] (vote) at (3.0,0)
    {\iflangja{多数決}{vote}};
  \draw[->] (b1.east) -- (b1.east -| vote.west);
  \draw[->] (b2.east) -- (vote.west);
  \draw[->] (b3.east) -- (b3.east -| vote.west);
  \draw[->] (vote.east) -- (4.1,0) node[right, lab] {\iflangja{出力}{out}};
  \end{scope}
\end{tikzpicture}

  \caption{$N$ モジュール冗長．(a) 2つのモジュールと比較器は，一方の
    モジュールにフォールトがあることは検出できるが，どちらであるかは
    分からない．(b) 3つのモジュールと多数決器は，いずれか1つのモジュールの
    誤った出力を隠蔽する．}
  \label{fig:l17-nmr}
\end{figure}

モジュールが2つの場合（\emph{2 重モジュール冗長}），不一致はどちらか一方に
フォールトがあることを示すが，どちらであるかは示さないので，エラーは検出
できても比較だけでは訂正できない．チェックポインティングと組み合わせれば，
システムは巻き戻して計算をやり直す．モジュールが3つの場合
（\emph{3 重モジュール冗長}），正しい結果は少なくとも2つのモジュールが
出した結果であり，フォールトのあるモジュールが1つであれば，巻き戻しなしに
多数決で退けられる---そのフォールトが永久的であってもである．\caution{$N$ 個のモジュールは1つの
設計を共有するので，設計フォールトは許容できない．}%
一般に，
$N$ 個のモジュールの多数決は最大 $\lfloor (N-1)/2 \rfloor$ 個のフォールト
のあるモジュールを隠蔽する（\cref{tab:l17-nmr}）．その代償は，ハードウェア
とエネルギーが $N$ 倍になることである．

\begin{table}[htbp]
  \centering\small
  \caption{$N$ モジュール冗長が扱えるフォールトのあるモジュール．}
  \label{tab:l17-nmr}
  \begin{tabular}{@{}cp{0.75\linewidth}@{}}
    \toprule
    $N$ & フォールトのあるモジュールの影響 \\
    \midrule
    2 & フォールトのあるモジュール 1 個を検出できる（出力が食い違う）が訂正はできない \\
    3 & フォールトのあるモジュール 1 個は 2 対 1 で退けられる \\
    5 & フォールトのあるモジュール 1 個は 4 対 1，2 個は 3 対 2 で退けられる \\
    \bottomrule
  \end{tabular}
\end{table}

フォールトのあるモジュールが1個であれば，5 モジュールのシステムはさらに
もう1個のフォールトを隠蔽できる．2個の場合，正しい3個がなお 3 対 2 で
多数を占めるが，余裕は残っていない．3個目にフォールトが生じれば，正しい
2個が退けられ，多数決は誤った結果を出してしまう．そのため，そうなる前に
運用を停止する．5 台の飛行制御計算機を備えたスペースシャトルはこの方針に
従っていた．1 台が故障しても任務は通常どおり続行され，2 台目が故障すると
任務は打ち切られた．\sidenote{同一の IBM AP-101 5 台のうち，同じプログラム
で多数決を行ったのは 4 台である．5 台目は，共有された設計フォールトに備えて
独立に書かれたソフトウェアを実行していた．}

\section{メモリとストレージのための符号}
\label{sec:l17-codes}

メモリとストレージについては，モジュール全体を複製するのは過剰である．
\source{Lesson 17，動画 14；H\&P 第2章．}フォールトは通常，格納された
ビットのうち数ビットを壊すだけであり，それらは誤り検出符号や誤り訂正符号
により，はるかに少ない追加ビットで保護できる．

\begin{definition}[パリティ]
  データビットの組 $d_1,\dots,d_n$ の\term[ぱりてぃ]{パリティ}[parity]と
  は，それらの排他的論理和
  \begin{equation}
    p \;=\; d_1 \oplus d_2 \oplus \dots \oplus d_n ,
    \label{eq:l17-parity}
  \end{equation}
  であり，パリティビットとは，$p$ をデータとともに格納する1つの追加ビット
  である．
\end{definition}

データを読み出すときにはパリティを再計算し，格納されているビットと比較
する．パリティビット自身を含め，どの1ビットが反転しても結果が変わるので，
1ビットの誤りはすべて検出される．パリティはどのビットが反転したかを示さ
ないので誤りを訂正することはできず，2ビットが反転した場合は互いに打ち消し
合って気づかれない．例えばバイト \texttt{10110010} は 1 が 4 個でパリティ
は 0 である．これが \texttt{10100010} として読み出されると，再計算した
パリティは 1 となり，誤りが検出される．

\begin{definition}[誤り訂正符号]
  \term[あやまりていせいふごう]{誤り訂正符号}[error-correcting code]%
  \index{ECC|see{誤り訂正符号}}（ECC）は，データビットの組に十分な数の
  検査ビットを付加し，限られた数の反転したビットの位置を特定して訂正
  できるようにしたものである．
\end{definition}

ECC 付きのメモリモジュールは通常 SECDED 符号（1 ビット訂正・2 ビット検出）
を用いる．例えば 64 データビットごとに 8 検査ビットを付加し，語の中の
1 ビットの反転は訂正し，2 ビットの反転は検出する．\tip{検査ビットの数：
$2^r \ge k+r+1$ を満たす最小の $r$ が，$k$ データビット中の 1 ビット反転の
位置を特定する．さらに 1 ビットで 2 ビット反転を検出する．$k=64$ なら
8 ビット，すなわち 72 ビット語である．}ディスクは各セクタを，
Reed--Solomon 符号のようなはるかに強力な符号で保護する．これらの符号は
多数の反転したビット，特に記録面の欠陥による隣接ビットのバースト誤りを
訂正する．こうした符号はセクタ内の壊れたビットに対処するものであり，
ディスクが丸ごと動かなくなった場合には役に立たない．その場合に必要なのは，
ディスクをまたぐ冗長性である．

\section{独立ディスク冗長アレイ}
\label{sec:l17-raid}

複数のディスクを1つの集合にまとめ，単一のディスクの役割を果たさせることが
できる．その集合は，構成するどのディスクよりも大容量に，高速に，あるいは
高信頼にできる．\source{Lesson 17，動画 15；H\&P §D.2；
\cite{patterson1988}．}このような集合を独立ディスク冗長アレイ，すなわち
\term{RAID}[redundant array of independent disks]%
\index{どくりつでぃすくじょうちょうあれい@独立ディスク冗長アレイ|see{RAID}}%
と呼ぶ．アレイには2つの目標がある．
\begin{itemize}
  \item 単一のディスクより高い性能を得ること．
  \item セクタが不良になっても，ディスクが丸ごと故障しても，読み書きの
    サービスを通常どおり提供すること．
\end{itemize}
すべての構成が両方の目標を追うわけではない．RAID 0%
（\cref{sec:l17-raid0}）は性能だけを目指す．

アレイ内の各ディスクは，依然として自身の符号でセクタを保護している
（\cref{sec:l17-codes}）．したがって，セクタが正しく読めないときや，
ディスクがまったく応答しないとき，アレイは\emph{どの}ディスクの\emph{どの}%
ブロックが影響を受けたのかを知っている．アレイの冗長性は誤りを見つける
必要はなく，位置の分かっているブロックの内容を復元しさえすればよい．
メモリのパリティビットが誤りを検出できるだけであるのに対し，アレイでは
1つのパリティブロックで誤りを訂正できるのは，このためである
（\cref{sec:l17-raid4}）．

RAID の各種の構成はレベルと呼ばれ，番号が付けられている．番号は構成を
識別するものであって，順位を表すものではない．この章ではレベル 0，1，4，
5，6 を扱う．\margin{\cite{patterson1988} では「I」は
\emph{inexpensive}（安価な）を表していた．レベル 2 と 3 はほとんど使われ
ないので，ここでは省く．}

\begin{keypoint}
  パリティビットは位置の分からない誤りを検出するだけだが，パリティブロック
  は位置の分かっている失われたデータを復元する．ディスクが自分の誤りを
  報告することが，パリティに基づく RAID を可能にしている．
\end{keypoint}

\section{RAID 0: ストライピング}
\label{sec:l17-raid0}

最初の構成は冗長性をまったく用いず，性能だけを目指す．\source{Lesson 17，
動画 16--18；H\&P §D.2．}1台の論理ディスクのデータを複数の物理ディスクに
分散させ，連続するブロックが異なるディスクに格納されるようにすることを
\term[すとらいぴんぐ]{ストライピング}[striping]という．全ディスクの同じ
位置にあるブロックは\emph{ストライプ}をなす．冗長性なしにストライピングを
用いる $N$ 台のディスクのアレイが \term{RAID 0}[RAID 0] である．論理
ブロック $i$ はディスク $i \bmod N$ に格納される
（\cref{fig:l17-raid01}a）．

\begin{figure}[htbp]
  \centering
  % Block placement in (a) RAID 0 with three disks and (b) RAID 1.
% Numbers are logical block numbers.
\begin{tikzpicture}[x=1cm,y=1cm, font=\footnotesize\sffamily]
  \def\blkw{0.95}\def\blkh{0.42}\def\blkgap{1.1}
  % ---------------- (a) RAID 0
  \node[anchor=west, font=\footnotesize\bfseries] at (-0.1,0.95)
    {\iflangja{(a) RAID 0（ストライピング）}{(a) RAID 0 (striping)}};
  \foreach \d in {0,1,2} {
    \node[font=\scriptsize, text=ln-muted] at ({\d*\blkgap+\blkw/2},0.3) {\iflangja{ディスク \d}{Disk \d}};
    \foreach \s in {0,...,3} {
      \pgfmathtruncatemacro{\b}{3*\s+\d}
      \draw[fill=white] ({\d*\blkgap},{-\s*\blkh}) rectangle ++(\blkw,-\blkh);
      \node at ({\d*\blkgap+\blkw/2},{-\s*\blkh-\blkh/2}) {\b};
    }
    \node[font=\scriptsize, text=ln-muted] at ({\d*\blkgap+\blkw/2},{-4*\blkh-0.15}) {$\vdots$};
  }
  % ---------------- (b) RAID 1
  \begin{scope}[xshift=5.6cm]
  \node[anchor=west, font=\footnotesize\bfseries] at (-0.1,0.95)
    {\iflangja{(b) RAID 1（ミラーリング）}{(b) RAID 1 (mirroring)}};
  \foreach \d in {0,1} {
    \node[font=\scriptsize, text=ln-muted] at ({\d*\blkgap+\blkw/2},0.3) {\iflangja{ディスク \d}{Disk \d}};
    \foreach \s in {0,...,3} {
      \draw[fill=white] ({\d*\blkgap},{-\s*\blkh}) rectangle ++(\blkw,-\blkh);
      \node at ({\d*\blkgap+\blkw/2},{-\s*\blkh-\blkh/2}) {\s};
    }
    \node[font=\scriptsize, text=ln-muted] at ({\d*\blkgap+\blkw/2},{-4*\blkh-0.15}) {$\vdots$};
  }
  \end{scope}
\end{tikzpicture}

  \caption{論理ブロックの配置．(a) 3 台のディスクの RAID 0．各ストライプは
    異なる3つのブロックを保持する．(b) RAID 1．両方のディスクがすべての
    ブロックを保持する．}
  \label{fig:l17-raid01}
\end{figure}

\subsection{性能}

アレイの容量は 1 台のディスクの $N$ 倍である．\margin{ストライピングの単位
は数十から数百キロバイトの\emph{チャンク}であり，1 セクタではない．図中の
各ブロックは1つのチャンクを表している．}異なるディスクにある
ブロックへの要求は同時に処理され，大きな転送は並行して進む $N$ 個の転送に
分割されるので，スループットは 1 台のディスクの最大 $N$ 倍になる．単一の
ブロックへのアクセスには，そのディスクでのシークと回転待ち（第16章）が
依然として必要である．したがって小さな要求が速く完了するのは，他の要求の
後ろで待ち行列に並ぶ時間が短くなるからにすぎない．大きな転送も，データが
$N$ 台のディスクとの間で同時に移動するので速く完了する．

\subsection{信頼性}

信頼性の計算では，ディスクは互いに独立に故障し，各ディスクは一定の
\emph{故障率} $f_1 = 1/\mathrm{MTTF}_1$ を持つと仮定する．ここで
$\mathrm{MTTF}_1$ は 1 台のディスクの MTTF である．\margin{実際のディスク
は寿命の初期と末期に故障しやすい．一定の故障率は，使用期間の長い中間部分を
モデル化している．}このとき，ディスクの集合の故障率は個々の故障率の和と
なる．RAID 0 では各ディスクが論理ディスクの $N$ ブロックごとに 1 ブロックを
保持するので，どのディスクが故障してもデータが失われる．アレイは $N f_1$
の率で故障するので，その MTTF は次式となる．
\begin{equation}
  \mathrm{MTTF}_{\text{RAID 0}} \;=\; \frac{\mathrm{MTTF}_1}{N} .
  \label{eq:l17-raid0}
\end{equation}
アレイにとって故障とはデータが失われることなので，その MTTF は
\emph{平均データ喪失時間}（MTTDL）とも呼ばれる．

\begin{example}
  \label{ex:l17-raid0}
  $\mathrm{MTTF}_1 = \SI{120000}{\hour}$（約 13.7 年）のディスク 5 台で
  RAID 0 のアレイを構成する．アレイは 1 台のディスクの最大 5 倍の速度と
  5 倍の容量を持つが，その MTTF は
  $\SI{120000}{\hour}/5 = \SI{24000}{\hour}$，約 2.7 年にすぎない．
\end{example}

\section{RAID 1: ミラーリング}
\label{sec:l17-raid1}

2つ目の構成は，ディスクを信頼性のために費やす．\source{Lesson 17，
動画 19--22；H\&P §D.2．}同じデータを2台のディスクに書き，各ディスクが
すべてのブロックの完全な複製を保持するようにすることを
\term[みらーりんぐ]{ミラーリング}[mirroring]という．ミラー化された2台の
ディスクのアレイ（\cref{fig:l17-raid01}b）は，いずれか1台が動作している
限りすべてのデータを利用可能に保つ．これを \term{RAID 1}[RAID 1] という．

\begin{itemize}
  \item \textbf{書き込み}は両方のディスクに対して行われる．2つの書き込みは
    並行して進むので，書き込みの速さは単一のディスクと同じである．
  \item \textbf{読み出し}はどちらの複製からでも処理できるので，読み出し
    スループットは 1 台のディスクの 2 倍である．
  \item \textbf{フォールト}は，1 台のディスクに限られていれば，そのディスク
    全体の故障も含めて許容される．複製が2つしかないので多数決はできないが，
    その必要もない．ディスク自身の符号がどちらの複製が不良かを報告するので，
    もう一方の複製を使えばよい．
  \item \textbf{コスト}: 2 台のディスクで 1 台分の容量が得られる．
\end{itemize}

\subsection{修復なしの信頼性}

両方のディスクが動作している間，アレイは $2f_1$ の率でディスク故障に
見舞われるので，最初の故障は平均して $\mathrm{MTTF}_1/2$ 後に起きる．
アレイは残り 1 台でこれを生き延びる．故障率は一定なので，生き残った
ディスクは新品同様であり，さらに平均 $\mathrm{MTTF}_1$ 後に故障する．
データが失われるのはそのときである．
\begin{equation}
  \mathrm{MTTF}_{\text{RAID 1}} \;=\; \frac{\mathrm{MTTF}_1}{2} + \mathrm{MTTF}_1
  \;=\; 1.5\,\mathrm{MTTF}_1 .
  \label{eq:l17-raid1-norepair}
\end{equation}

\subsection{修復ありの信頼性}

実際には，故障したディスクは交換され，生き残ったディスクから新しい
ディスクへデータが複写される．故障から，新しいディスクが完全な複製を保持
するまでの平均時間を $\mathrm{MTTR}_1$ とする．\margin{$\mathrm{MTTR}_1$ は
保守員を呼ぶ時間ではなく複写の時間である．16~TB のディスクを 150~MB/s で
埋め直せば約 \SI{30}{\hour} かかる．以下の例で用いる値はこれである．}%
データが失われるのは，この
期間内に2台目のディスクが故障した場合だけである（$N=2$ とした
\cref{fig:l17-raid-repair}）．$\mathrm{MTTR}_1 \ll \mathrm{MTTF}_1$ なので
その確率はおよそ $\mathrm{MTTR}_1/\mathrm{MTTF}_1$ であり，平均すると
$\mathrm{MTTF}_1/\mathrm{MTTR}_1$ 回の最初の故障が起きるうちの1回が
データ喪失に至る．そのそれぞれの前には，両方のディスクが動作している
$\mathrm{MTTF}_1/2$ の期間がある．したがって次式が得られる．
\begin{equation}
  \mathrm{MTTF}_{\text{RAID 1}} \;\approx\;
  \frac{\mathrm{MTTF}_1}{2} \times \frac{\mathrm{MTTF}_1}{\mathrm{MTTR}_1}
  \;=\; \frac{\mathrm{MTTF}_1^2}{2\,\mathrm{MTTR}_1} .
  \label{eq:l17-raid1-repair}
\end{equation}

\begin{figure}[htbp]
  \centering
  % States of an array that tolerates one disk failure (RAID 1 with N = 2,
% RAID 4/5 with N disks) when failed disks are replaced.
\begin{tikzpicture}[x=1cm,y=1cm, font=\footnotesize\sffamily, >={Stealth[length=1.8mm]},
  st/.style={draw, rounded corners=3pt, minimum height=0.9cm, text width=2.2cm,
    align=center, inner sep=3pt},
  lab/.style={font=\scriptsize, align=center}]
  \node[st, fill=ln-accent!15] (ok) at (0,0) {\iflangja{$N$ 台すべて動作}{all $N$ disks\\working}};
  \node[st, fill=ln-box] (rep) at (5.4,0) {\iflangja{1 台が故障\\（修復中）}{one disk failed\\(under repair)}};
  \node[st, fill=ln-caution!15] (lost) at (5.4,-2.6) {\iflangja{データ喪失}{data lost}};
  \draw[->] (ok) to[bend left=18] node[lab, above]
    {\iflangja{$N$ 台のいずれかが故障\\（平均 $\mathrm{MTTF}_1/N$ 後）}{one of the $N$ disks fails\\(after $\mathrm{MTTF}_1/N$ on average)}} (rep);
  \draw[->] (rep) to[bend left=18] node[lab, below]
    {\iflangja{修復完了\\（$\mathrm{MTTR}_1$ 後）}{repair completes\\(after $\mathrm{MTTR}_1$)}} (ok);
  \draw[->, ln-caution] (rep) -- node[lab, right, text=black]
    {\iflangja{修復完了前に残り $N-1$ 台の\\いずれかが故障（確率\\$\approx (N-1)\,\mathrm{MTTR}_1/\mathrm{MTTF}_1$）}{one of the other $N-1$ disks\\fails before the repair completes\\(probability $\approx (N-1)\,\mathrm{MTTR}_1/\mathrm{MTTF}_1$)}} (lost);
\end{tikzpicture}

  \caption{1 台のディスク故障を許容するアレイと，その修復．データが失われる
    のは，1 台目を修復している間に 2 台目が故障した場合だけである．RAID 1
    では $N=2$ である．}
  \label{fig:l17-raid-repair}
\end{figure}

\begin{example}
  \label{ex:l17-raid1}
  \cref{ex:l17-raid0} のディスク（$\mathrm{MTTF}_1 = \SI{120000}{\hour}$）%
  \margin{データシートの MTTF \num{1000000}--\num{1500000}~h は年間故障率
  \SI{1}{\percent} 未満を意味する．実地で測定された交換率は
  \SIrange{2}{4}{\percent} であり，MTTF はむしろ \num{300000}~h に近い
  \cite{schroeder2007}．}%
  を用いると，修復なしの RAID 1 アレイの MTTF は \SI{180000}{\hour}，
  約 20.5 年である．故障したディスクが交換され，その複製が
  $\mathrm{MTTR}_1 = \SI{30}{\hour}$ 以内に復元されるなら，MTTF は
  $\SI{60000}{\hour} \times 4000 = \SI{2.4e8}{\hour}$，約 \num{27000} 年に
  なる．
\end{example}

\needspace{9\baselineskip}
\begin{keypoint}
  修復がなければ冗長性はほとんど役に立たない．ミラー化された2台のディスク
  は平均して $1.5\,\mathrm{MTTF}_1$ しかもたない．そのうち1台目が
  $\mathrm{MTTF}_1/2$ 後に故障するからである．修復があれば，冗長性は MTTF
  を $\mathrm{MTTF}_1/\mathrm{MTTR}_1$ 程度の倍率で引き上げる．フォールト
  トレラントなシステムでは，モジュールの修復を速くすることが，システムの
  可用性だけでなく信頼性も改善する．
\end{keypoint}

\section{RAID 4: ブロックインターリーブパリティ}
\label{sec:l17-raid4}

パリティは，ミラーリングと同じ冗長性をその何分の一かのコストで提供する．
\source{Lesson 17，動画 23--27；H\&P §D.2．}$N$ 台のディスクのアレイでは，
そのうち $N-1$ 台を RAID 0 と同様にストライピングしたデータディスクとし，
残る 1 台をパリティディスクとすることができる．パリティディスクの各
ストライプのブロックは，そのストライプのデータブロックの排他的論理和で
ある（\cref{fig:l17-raid4-write}a）．この構成を \term{RAID 4}[RAID 4] と
いう．データブロック $D_0,\dots,D_{N-2}$ からなるストライプに
\cref{eq:l17-parity} をビットごとに適用すると次式が得られる．
\begin{equation}
  P \;=\; D_0 \oplus D_1 \oplus \dots \oplus D_{N-2} .
  \label{eq:l17-raid4-parity}
\end{equation}

$x \oplus x = 0$ であるから，\cref{eq:l17-raid4-parity} の両辺に残存する
すべてのブロックの排他的論理和を取ると，失われたブロックが，パリティを
含む同じストライプの他のすべてのブロックの排他的論理和に等しいことが
分かる．\caution{符号理論では，位置の分かっている喪失を\emph{消失}
（erasure）と呼ぶ．検査シンボル1個で消失は1個訂正できるが，位置の分からない
誤りを1個，位置の特定まで含めて訂正するには2個必要である．}%
ディスクが故障しても，その上のすべてのブロックは他の $N-1$ 台の
ディスクから復元できる．故障したのがパリティディスクであれば，パリティを
計算し直すだけでよい．冗長性のコストは $N$ 台のうち 1 台であり，ミラー
リングの2つ目の複製よりはるかに小さい．

\subsection{書き込み}

書き込みは，そのストライプのパリティを更新しなければならない．ストライプの
すべてのデータブロックからパリティを計算し直すには，他の $N-2$ 台の
ディスクを読む必要がある．そうする代わりに，新しいパリティは，古いパリティ
から古いデータを取り除き，新しいデータを加えることで得られる
（\cref{fig:l17-raid4-write}b）．
\begin{equation}
  P' \;=\; P \oplus D \oplus D' .
  \label{eq:l17-parity-update}
\end{equation}
したがって書き込みは，古いデータブロックと古いパリティブロックを読み，
新しいデータブロックと新しいパリティブロックを書く．$N$ によらず，2 回の
読み出しと 2 回の書き込みである．\margin{この読み出し・変更・書き込みの
手順は\emph{スモールライト問題}として知られる．ストライプ全体を覆う書き込み
はこれを免れる．新しいパリティが新しいデータだけから求まるからである．}

\begin{figure}[htbp]
  \centering
  % (a) RAID 4 with three data disks and a parity disk; rows are stripes A-D.
% (b) A small write to block B1: two reads, an XOR, two writes.
\begin{tikzpicture}[x=1cm,y=1cm, font=\footnotesize\sffamily, >={Stealth[length=1.8mm]},
  blk/.style={draw, fill=white, minimum width=1.25cm, minimum height=0.5cm, inner sep=1pt},
  lab/.style={font=\scriptsize, text=ln-muted, align=center}]
  \def\blkw{0.8}\def\blkh{0.45}\def\blkgap{0.9}
  % ---------------- (a) layout
  \node[anchor=west, font=\footnotesize\bfseries] at (-0.9,0.95) {(a) \iflangja{配置}{Layout}};
  \node[font=\scriptsize, text=ln-muted, anchor=east] at (-0.05,0.3) {\iflangja{ディスク}{Disk}};
  \foreach \d in {0,1,2,3} {
    \node[font=\scriptsize, text=ln-muted] at ({\d*\blkgap+\blkw/2},0.3) {\d};
  }
  \foreach \S [count=\r from 0] in {A,B,C,D} {
    \node[font=\scriptsize, text=ln-muted, anchor=east] at (-0.05,{-\r*\blkh-\blkh/2}) {\S};
    \foreach \d in {0,1,2} {
      \draw[fill=white] ({\d*\blkgap},{-\r*\blkh}) rectangle ++(\blkw,-\blkh);
      \node at ({\d*\blkgap+\blkw/2},{-\r*\blkh-\blkh/2}) {$\S_\d$};
    }
    \draw[fill=ln-accent!15] ({3*\blkgap},{-\r*\blkh}) rectangle ++(\blkw,-\blkh);
    \node at ({3*\blkgap+\blkw/2},{-\r*\blkh-\blkh/2}) {$P_\S$};
  }
  \draw[ln-caution, very thick] ({1*\blkgap},{-\blkh}) rectangle ++(\blkw,-\blkh);
  \draw[ln-caution, very thick] ({3*\blkgap},{-\blkh}) rectangle ++(\blkw,-\blkh);
  \node[lab] at ({1*\blkgap+\blkw/2},{-4*\blkh-0.3}) {\iflangja{データ}{data}};
  \node[lab] at ({3*\blkgap+\blkw/2},{-4*\blkh-0.3}) {\iflangja{パリティ}{parity}};
  % ---------------- (b) small write
  \begin{scope}[xshift=5.4cm]
  \node[anchor=west, font=\footnotesize\bfseries] at (-0.4,0.95) {(b) \iflangja{$B_1$ への書き込み}{Write to $B_1$}};
  \node[blk] (nd) at (1.6,0.1)  {\iflangja{新 $B_1'$}{new $B_1'$}};
  \node[blk] (od) at (1.6,-0.7) {\iflangja{旧 $B_1$}{old $B_1$}};
  \node[blk] (op) at (1.6,-1.5) {\iflangja{旧 $P_B$}{old $P_B$}};
  \node[lab, anchor=east] at (0.9,0.1)  {\iflangja{ディスク 1 に\\書く}{written\\to disk 1}};
  \node[lab, anchor=east] at (0.9,-0.7) {\iflangja{ディスク 1 から\\読む}{read from\\disk 1}};
  \node[lab, anchor=east] at (0.9,-1.5) {\iflangja{ディスク 3 から\\読む}{read from\\disk 3}};
  \node[draw, circle, fill=ln-accent!15, inner sep=0pt, minimum size=0.55cm] (x) at (3.3,-0.7) {$\oplus$};
  \draw[->] (nd.east) -- (x);
  \draw[->] (od.east) -- (x);
  \draw[->] (op.east) -- (x);
  \node[blk, fill=ln-accent!15] (np) at (5.1,-0.7) {\iflangja{新 $P_B'$}{new $P_B'$}};
  \draw[->] (x) -- (np);
  \node[lab, anchor=north] at (5.1,-1.0) {\iflangja{ディスク 3 に\\書く}{written\\to disk 3}};
  \end{scope}
\end{tikzpicture}

  \caption{RAID 4．(a) 3 台のデータディスクと 1 台のパリティディスク．行が
    ストライプである．(b) ブロック $B_1$ への書き込みは，$B_1$ と $P_B$ を
    読み，\cref{eq:l17-parity-update} で新しいパリティを計算し，$B_1'$ と
    $P_B'$ を書く．}
  \label{fig:l17-raid4-write}
\end{figure}

\begin{example}
  \label{ex:l17-parity}
  3 台のデータディスクと 1 台のパリティディスクからなるアレイと，4 ビットの
  ブロックからなる1つのストライプを考える（\cref{tab:l17-parity}）．
  パリティは $1001 \oplus 0111 \oplus 1010 = 0100$ である．ディスク 1 が
  故障すると，そのブロックは $1001 \oplus 1010 \oplus 0100 = 0111$ として
  復元される．ディスク 1 に $0010$ を書き込むと，
  \cref{eq:l17-parity-update} により新しいパリティは
  $0100 \oplus 0111 \oplus 0010 = 0001$ となり，これは新しいデータ
  ブロックの排他的論理和 $1001 \oplus 0010 \oplus 1010$ に等しい．
\end{example}

\begin{table}[htbp]
  \centering\footnotesize
  \caption{\cref{ex:l17-parity} における RAID 4 の1つのストライプ．復元
    および更新されたブロックを太字で示す．}
  \label{tab:l17-parity}
  \begin{tabular}{@{}lcccc@{}}
    \toprule
    & ディスク 0 & ディスク 1 & ディスク 2 & \shortstack{ディスク 3\\（パリティ）} \\
    \midrule
    初期のストライプ                  & 1001 & 0111 & 1010 & 0100 \\
    ディスク 1 の喪失と復元           & 1001 & \textbf{0111} & 1010 & 0100 \\
    ディスク 1 に 0010 を書き込んだ後 & 1001 & \textbf{0010} & 1010 & \textbf{0001} \\
    \bottomrule
  \end{tabular}
\end{table}

\subsection{性能と信頼性}

\begin{itemize}
  \item \textbf{読み出し}はデータディスクだけを使うので，読み出し
    スループットは 1 台のディスクの $N-1$ 倍である．
  \item \textbf{書き込み}はすべてパリティディスクを必要とし，読み出しに
    1 回，書き込みに 1 回アクセスする．異なるデータディスクへの書き込みで
    あっても並行しては進められず，各書き込みがパリティディスクを 2 回の
    アクセスのあいだ占有する．したがって書き込みスループットは，$N$ の値に
    よらず単一のディスクの半分である．
  \item \textbf{信頼性}: アレイはどれか 1 台のディスクの故障を許容する．
    修復がなければ，最初の故障は平均 $\mathrm{MTTF}_1/N$ 後に，2 回目は
    さらに $\mathrm{MTTF}_1/(N-1)$ 後に起きる．$N \ge 3$ ではこの和が
    $\mathrm{MTTF}_1$ より小さい．すなわち修復のないこのようなアレイは，
    単一のディスクより信頼性が低い．修復があれば，
    \cref{eq:l17-raid1-repair} の議論が，最初に故障しうるディスクが $N$ 台，
    修復中に故障しうるディスクが $N-1$ 台という形で適用でき，次式が得られる
    （\cref{fig:l17-raid-repair}）．
    \begin{equation}
      \mathrm{MTTF}_{\text{RAID 4}} \;\approx\;
      \frac{\mathrm{MTTF}_1}{N} \times
      \frac{\mathrm{MTTF}_1}{(N-1)\,\mathrm{MTTR}_1} .
      \label{eq:l17-raid4}
    \end{equation}
\end{itemize}

\section{RAID 5: 分散パリティ}
\label{sec:l17-raid5}

RAID 4 のパリティディスクは書き込みのボトルネックである．\source{Lesson
17，動画 28--29；H\&P §D.2．}その対策は，パリティを全ディスクに分散させる
ことである．各ストライプが1つのパリティブロックを持つことは変わらず，
その計算も RAID 4 とまったく同じだが，異なるストライプのパリティブロックは
異なるディスクに格納され，$N$ 台のディスクを巡回する
（\cref{fig:l17-raid5-6}a）．この構成を \term{RAID 5}[RAID 5] という．
どのディスクもデータとパリティの両方を保持する．

\begin{figure}[htbp]
  \centering
  % (a) RAID 5 with five disks: the parity block rotates among the disks; two
% writes to different stripes use disjoint disks.  (b) RAID 6 with two check
% blocks (P and Q) per stripe.
\begin{tikzpicture}[x=1cm,y=1cm, font=\footnotesize\sffamily]
  \def\blkw{0.86}\def\blkh{0.45}\def\blkgap{0.94}
  % ---------------- (a) RAID 5
  \node[anchor=west, font=\footnotesize\bfseries] at (-0.5,0.95) {(a) RAID 5};
  \node[font=\scriptsize, text=ln-muted, anchor=east] at (-0.05,0.3) {\iflangja{ディスク}{Disk}};
  \foreach \d in {0,...,4} {
    \node[font=\scriptsize, text=ln-muted] at ({\d*\blkgap+\blkw/2},0.3) {\d};
  }
  % stripe letter / parity disk
  \foreach \S/\p [count=\r from 0] in {A/4,B/3,C/2,D/1,E/0} {
    \node[font=\scriptsize, text=ln-muted, anchor=east] at (-0.05,{-\r*\blkh-\blkh/2}) {\S};
    \foreach \d in {0,...,4} {
      \pgfmathtruncatemacro{\i}{\d<\p ? \d : \d-1}
      \ifnum\d=\p
        \draw[fill=ln-accent!15] ({\d*\blkgap},{-\r*\blkh}) rectangle ++(\blkw,-\blkh);
        \node at ({\d*\blkgap+\blkw/2},{-\r*\blkh-\blkh/2}) {$P_\S$};
      \else
        \draw[fill=white] ({\d*\blkgap},{-\r*\blkh}) rectangle ++(\blkw,-\blkh);
        \node at ({\d*\blkgap+\blkw/2},{-\r*\blkh-\blkh/2}) {$\S_\i$};
      \fi
    }
  }
  % write to A0 (disks 0, 4) and to C1 (disks 1, 2)
  \draw[ln-caution, very thick] (0,0) rectangle ++(\blkw,-\blkh);
  \draw[ln-caution, very thick] ({4*\blkgap},0) rectangle ++(\blkw,-\blkh);
  \draw[ln-tip, very thick] ({1*\blkgap},{-2*\blkh}) rectangle ++(\blkw,-\blkh);
  \draw[ln-tip, very thick] ({2*\blkgap},{-2*\blkh}) rectangle ++(\blkw,-\blkh);
  % ---------------- (b) RAID 6
  \begin{scope}[xshift=6.3cm]
  \node[anchor=west, font=\footnotesize\bfseries] at (-0.5,0.95) {(b) RAID 6};
  \node[font=\scriptsize, text=ln-muted, anchor=east] at (-0.05,0.3) {\iflangja{ディスク}{Disk}};
  \foreach \d in {0,...,4} {
    \node[font=\scriptsize, text=ln-muted] at ({\d*\blkgap+\blkw/2},0.3) {\d};
  }
  % stripe letter / P disk / Q disk
  \foreach \S/\p/\q [count=\r from 0] in {A/3/4,B/2/3,C/1/2,D/0/1,E/4/0} {
    \node[font=\scriptsize, text=ln-muted, anchor=east] at (-0.05,{-\r*\blkh-\blkh/2}) {\S};
    \foreach \d in {0,...,4} {
      \pgfmathtruncatemacro{\i}{(\d>\p ? 1 : 0) + (\d>\q ? 1 : 0)}
      \pgfmathtruncatemacro{\i}{\d-\i}
      \ifnum\d=\p
        \draw[fill=ln-accent!15] ({\d*\blkgap},{-\r*\blkh}) rectangle ++(\blkw,-\blkh);
        \node at ({\d*\blkgap+\blkw/2},{-\r*\blkh-\blkh/2}) {$P_\S$};
      \else\ifnum\d=\q
        \draw[fill=ln-accent!30] ({\d*\blkgap},{-\r*\blkh}) rectangle ++(\blkw,-\blkh);
        \node at ({\d*\blkgap+\blkw/2},{-\r*\blkh-\blkh/2}) {$Q_\S$};
      \else
        \draw[fill=white] ({\d*\blkgap},{-\r*\blkh}) rectangle ++(\blkw,-\blkh);
        \node at ({\d*\blkgap+\blkw/2},{-\r*\blkh-\blkh/2}) {$\S_\i$};
      \fi\fi
    }
  }
  \end{scope}
\end{tikzpicture}

  \caption{(a) 5 台のディスクの RAID 5．パリティブロックはディスクを巡回
    する．$A_0$ への書き込み（赤）と $C_1$ への書き込み（緑）は使う
    ディスクが重ならないので，並行して進められる．(b) ストライプごとに2つの
    チェックブロック $P$ と $Q$ を持つ RAID 6．}
  \label{fig:l17-raid5-6}
\end{figure}

\begin{itemize}
  \item \textbf{読み出し}: データは $N$ 台すべてのディスクに分散している
    ので，読み出しスループットは 1 台のディスクの $N$ 倍である．
  \item \textbf{書き込み}には依然として 2 回の読み出しと 2 回の書き込みが
    必要だが（\cref{eq:l17-parity-update}），パリティを保持するディスクは
    ストライプによって異なる．異なるストライプへの書き込みは通常，異なる
    ディスクを使うので並行して進む．書き込み 1 回当たりの 4 回のアクセスが
    $N$ 台すべてのディスクに分散するので，書き込みスループットは 1 台の
    ディスクの $N/4$ 倍である．
  \item \textbf{信頼性}と\textbf{コスト}は RAID 4 と同じである．1 台の
    ディスク故障を許容し，修復ありの MTTF は \cref{eq:l17-raid4} で与えられ，
    ディスク 1 台分の容量がパリティに使われる．
\end{itemize}

\begin{example}
  \label{ex:l17-raid5}
  6 台のディスクがあり，各ディスクは毎秒 200 回のアクセスを処理でき，
  $\mathrm{MTTF}_1 = \SI{120000}{\hour}$，$\mathrm{MTTR}_1 =
  \SI{30}{\hour}$ である．RAID 4 として構成すると毎秒 1000 回の読み出し
  または 100 回の書き込みを処理し，RAID 5 として構成すると毎秒 1200 回の
  読み出しまたは 300 回の書き込みを処理する．いずれの場合も修復ありの MTTF
  は $(\num{120000}/6) \times \num{120000}/(5 \times 30) =
  \SI{1.6e7}{\hour}$，約 \num{1800} 年である．
\end{example}

\begin{keypoint}
  RAID 5 は RAID 4 と同じ信頼性とコストを持ち，読み出しと書き込みの性能は
  より高い．あらゆる点で RAID 4 に劣らないので，RAID 4 が使われることは
  まれである．
\end{keypoint}

\section{RAID 6: 2つのチェックブロック}
\label{sec:l17-raid6}

RAID 5 は 1 台のディスク故障を許容する．次の構成は 2 台の故障を許容する．
\source{Lesson 17，動画 30--32；H\&P §D.2．}各ストライプに2つ目のチェック
ブロックを加えるのである（\cref{fig:l17-raid5-6}b）．一方のチェック
ブロック $P$ はパリティであり，もう一方の $Q$ は別の符号で計算される．この
構成を \term{RAID 6}[RAID 6] という．$P$ と $Q$ はストライプごとに2つの
独立な方程式を与える\margin{$Q$ は通常 $\mathrm{GF}(2^8)$ 上の
Reed--Solomon 検査であり，固定の $g$ を用いて $Q = \sum_i g^i D_i$ とする．
重みが異なることが2つ目の方程式を独立にしている．2つ目も排他的論理和では，
1つ目の繰り返しにしかならない．}ので，1つのストライプのどの2ブロックが
失われても復元できる．1ブロックだけが失われた場合は RAID 5 と同様にパリティから復元し，
2ブロックが失われた場合は両方の方程式を解いて求める．

\begin{itemize}
  \item \textbf{信頼性}: アレイはどの 2 台のディスクの故障も許容する．
  \item \textbf{コスト}: ディスク 2 台分の容量がチェックブロックに使われる．
  \item \textbf{書き込み}は両方のチェックブロックを更新しなければならない．
    RAID 5 では読み出しも書き込みも 2 回だったのに対し，3 回の読み出し
    （古いデータ，$P$，$Q$）と 3 回の書き込みが必要である．書き込み 1 回
    当たり 6 回のアクセスが $N$ 台のディスクに分散するので，書き込み
    スループットは 1 台のディスクのおよそ $N/6$ 倍となる．読み出しは
    $N$ 倍のままである．\tip{\cref{tab:l17-raid} の書き込みの列は，$N$ を
    書き込み1回当たりのアクセス数で割った値である．RAID 4 と RAID 5 は
    4 回，RAID 6 は 6 回である．RAID 4 が代わりに $1/2$ なのは，4 回すべてが
    1 台のパリティディスクに集中するからである．}
\end{itemize}

\subsection{RAID 6 と相関のある故障}

独立なディスク故障に対しては，修復ありの RAID 5 でもすでに千年を優に超える
MTTF が得られる（\cref{ex:l17-raid5}）．RAID 6 では，2 台を修復している間に
3 台目が故障した場合にのみデータが失われるので，MTTF はさらに
$\mathrm{MTTF}_1/((N-2)\,\mathrm{MTTR}_1)$ 程度の倍率で伸びる．例の数値では
約 1000 倍である．独立な故障に対しては，RAID 6 は過剰である．

RAID 6 の価値は，最初の故障と独立\emph{でない}故障にある．\margin{修復では
残りのディスクを端から端まで読むことにもなる．コンシューマ向けのデータ
シートが許す回復不能読み出し誤りは $10^{14}$ ビットに 1 回であり，7 台の
4~TB ディスクはその 2 倍のビット数を保持する．RAID 5 の再構築そのものが
ブロックを失いうる．}典型的なのは
修復中の運用フォールトで，運用者が，故障したディスクではなく動作している
ディスクを交換してしまう場合である．RAID 5 のアレイはこれで 2 台を失い，
データも失われる．RAID 6 のアレイは動作し続け，なお修復できる．

\begin{table}[htbp]
  \centering\small
  \caption{$\mathrm{MTTF}_1$ と $\mathrm{MTTR}_1$ を持つ $N$ 台のディスクに
    よる RAID レベルの比較．容量とスループットは 1 台のディスクに対する
    相対値である．MTTF は RAID 1，4，5，6 については修復ありを仮定して
    いる．RAID 0 は最初の故障でデータを失う．}
  \label{tab:l17-raid}
  \setlength{\tabcolsep}{4pt}
  \begin{tabular}{@{}lccccc@{}}
    \toprule
    レベル & 容量 & 読み出し & 書き込み & 許容故障台数 & MTTF \\
    \midrule
    RAID 0 & $N$   & $N$   & $N$   & 0 & $\mathrm{MTTF}_1/N$ \\
    RAID 1 ($N=2$) & 1 & 2 & 1   & 1 & $\mathrm{MTTF}_1^2/(2\,\mathrm{MTTR}_1)$ \\
    RAID 4 & $N-1$ & $N-1$ & $1/2$ & 1 & $\mathrm{MTTF}_1^2/(N(N-1)\,\mathrm{MTTR}_1)$ \\
    RAID 5 & $N-1$ & $N$   & $N/4$ & 1 & $\mathrm{MTTF}_1^2/(N(N-1)\,\mathrm{MTTR}_1)$ \\
    RAID 6 & $N-2$ & $N$   & $N/6$ & 2 & RAID 5 よりはるかに高い \\
    \bottomrule
  \end{tabular}
\end{table}

\needspace{10\baselineskip}
\begin{keypoint}[章のまとめ]
  モジュールのフォールトは，活性化されるとエラーになり，エラーが提供
  サービスに到達すると故障になる．信頼性は MTTF で，可用性は
  $\mathrm{MTTF}/(\mathrm{MTTF}+\mathrm{MTTR})$ で測られる．フォールトは
  原因（ハードウェア，設計，運用，環境）と継続時間（永久，間欠，過渡）で
  分類され，回避・許容・修復の高速化によって対処される．チェック
  ポインティングと $N$ モジュール冗長は一般のフォールトを許容し，パリティ
  と ECC はメモリを保護し，RAID はディスクの故障から守る．ストライピング
  （RAID 0）は信頼性と引き換えに速度を得，ミラーリング（RAID 1）は2つ目の
  複製で信頼性を買い，パリティ（RAID 4，5）はディスク 1 台分のコストで
  それを実現する．RAID 5 はパリティディスクのボトルネックを回避する．修復
  があれば，冗長アレイの MTTF は $\mathrm{MTTF}_1/\mathrm{MTTR}_1$ に応じて
  伸びる（\cref{tab:l17-raid}）．RAID 6 は主に，修復中の誤りのような相関の
  ある故障から守る．
\end{keypoint}

\end{document}
